Die Durchführung des Praxisprojekts im Umfeld der snoopmedia GmbH ermöglichte eine intensive Anwendung und Vertiefung der im Studium der Wirtschaftsinformatik erworbenen Kompetenzen. 
Insbesondere die Kernbereiche der Softwareentwicklung und des IT-Projektmanagements standen hierbei im Fokus.

\subsection{Software Engineering und Architektur}
Den größten Bezugspunkt bildeten die Module \textit{Software Engineering 1 \& 2}. Die theoretischen Konzepte zur Softwarearchitektur waren essenziell für die Entwicklung der Game-Engine.
 Hierbei kam das Prinzip der \textit{Separation of Concerns} (Trennung der Belange) zur Anwendung: Die strikte Trennung zwischen dem Datenmodell im Backend (PHP/Symfony), der 
 Logikschicht in den TypeScript-Services und der Präsentationsschicht im HTML5-Canvas spiegelt moderne Architekturmuster wider, wie sie im Studium vermittelt wurden.

Besonders hervorzuheben sind die \textit{Architekturentscheidungen} im Bereich des Entwurfs:
\begin{itemize}
    \item \textit{Entwurfsmuster:} Bei der Implementierung der Game-Engine wurden gezielt Entwurfsmuster eingesetzt, die dem Industriestandard für wiederverwendbare Software entsprechen 
    \cite{Gamma1994}. Das eingesetzte \textit{State-Pattern} ermöglichte die Verwaltung der verschiedenen Spielzustände, während das \textit{Observer-Pattern} für die Event-Steuerung 
    innerhalb der Szenen genutzt wurde. Diese praktische Anwendung vertiefte das Verständnis aus der Vorlesung \textit{Programmierung 2}.
    \item \textit{Datenmodellierung:} Die Erweiterung des Datenbankschemas erforderte fundierte Kenntnisse aus dem Modul \textit{Datenbanken}. Die Normalisierung der Tabellen 
    \texttt{world}, \texttt{scene} und \texttt{content} sowie die Umsetzung komplexer Relationen mittels Doctrine ORM zeigten die Relevanz einer sauberen Datenmodellierung für 
    die Skalierbarkeit von Systemen.
    \item \textit{Typsicherheit:} Durch den Einsatz von TypeScript wurden Konzepte der statischen Typisierung und Schnittstellendefinition (\textit{Interfaces}) angewendet, um die 
    Fehleranfälligkeit bei der Kommunikation mit der API zu minimieren. Ein direkter Transfer aus den Programmierkursen des Basisstudiums.
\end{itemize}

\subsection{Agile Methoden und IT-Projektmanagement}
Neben der rein technischen Umsetzung spielte das \textit{Informationsmanagement} und die Organisation der Arbeit eine zentrale Rolle. In der Zusammenarbeit mit der snoopmedia GmbH 
fand das Wissen über \textit{agiles Arbeiten} direkte Anwendung.

Obwohl kein striktes Framework wie Scrum nach Lehrbuch verfolgt wurde, orientierte sich der Arbeitsalltag an agilen Prinzipien. Dies beinhaltete das \textit{Ticketbasiertes Arbeiten}.
Die Nutzung eines Ticketsystems zur Aufgabensteuerung erlaubte eine transparente Dokumentation des Fortschritts und eine effiziente 
Priorisierung von Features und Bugfixes.

\subsection{Reflexion}
Zusammenfassend lässt sich festhalten, dass das Praxisprojekt das im Studium theoretisch gelernte Wissen „greifbar“ gemacht hat. Die Herausforderung, eine alleinige technische 
Verantwortung für ein Teilmodul eines Forschungsprojekts zu tragen, förderte nicht nur das Verständnis für komplexe IT-Infrastrukturen, sondern auch die Fähigkeit, eigenständig 
fundierte technische Entscheidungen auf Basis wissenschaftlicher Methoden zu treffen.